% Part: normal-modal-logic
% Chapter: tableaux
% Section: simple-S5

\documentclass[../../../include/open-logic-section]{subfiles}

\begin{document}

\olfileid{nml}{tab}{s5}

\olsection{\Log{S5} 的简化真值树}

\Log{S5} 相对于普遍模型类——即任意世界都可从任意世界可达的模型——是可靠且完备的。
在普遍模型中，可达关系并不重要：“存在一个世界~$w$ 使得 $\mSat{M}{!A}[w]$”为真，当且仅当存在这样一个从~$u$ 可达的 $w$。
因此在 \Log{S5} 中，我们可以把模型简单地定义为一个世界集合和一个赋值~$V$。
这表明，我们也应该能够简化分析表规则。
在一般情形中，我们取正整数序列作为前缀，以便追踪哪些前缀命名了可从其他世界到达的世界：$\sigma.n$ 命名一个从~$\sigma$ 可达的世界。
但在 \Log{S5} 中，任意世界都可从任意世界可达，因此无需如此追踪。
我们可以直接使用正整数作为前缀。
简化后的规则见\olref{tab:rules-S5}。

\begin{table}
  \begin{center}
    \def\arraystretch{3}\def\fCenter{}
    \begin{tabular}{|c|c|}
    \hline
    \iftag{prvBox}{
      \AxiomC{\sFmla{\True}{\Box !A}[n]}
      \RightLabel{\TRule{\True}{\Box}}
      \UnaryInfC{\sFmla{\True}{!A}[m]}
      \DisplayProof
      &
      \AxiomC{\sFmla{\False}{\Box !A}[n]}
      \RightLabel{\TRule{\False}{\Box}}
      \UnaryInfC{\sFmla{\False}{!A}[m]}
      \DisplayProof\\
      $m$ 已使用 & $m$ 是新的
      \\[1ex]
      \hline}{}
    \iftag{prvDiamond}{
      \AxiomC{\sFmla{\True}{\Diamond !A}[n]}
      \RightLabel{\TRule{\True}{\Diamond}}
      \UnaryInfC{\sFmla{\True}{!A}[m]}
      \DisplayProof
      &
      \AxiomC{\sFmla{\False}{\Diamond !A}[n]}
      \RightLabel{\TRule{\False}{\Diamond}}
      \UnaryInfC{\sFmla{\False}{!A}[m]}
      \DisplayProof\\
      $m$ 是新的 & $m$ 已使用
      \\[1ex]
      \hline}{}
    \end{tabular}
  \end{center}
  \caption{\Log{S5} 的简化规则。}
  \ollabel{tab:rules-S5}
\end{table}

\begin{ex}
  我们给出一个简化的闭分析表，来证明 $\Log{S5} \Proves
  \Ax{5}$，即 $\Diamond!A \lif \Box\Diamond!A$。
  \begin{oltableau}
    [\pFmla{\False}{\Diamond\formula{A} \lif \Box\Diamond \formula{A}}{1},
      just = \TAss
      [\pFmla{\True}{\Diamond \formula{A}}{1}, just = {\TRule{\False}{\lif}[1]}
        [\pFmla{\False}{\Box\Diamond \formula{A}}{1},
          just = {\TRule{\False}{\lif}[1]}
          [\pFmla{\False}{\Diamond \formula{A}}{2},
            just = {\TRule{\False}{\Box}[3]}
            [\pFmla{\True}{\formula{A}}{3},
              just = {\TRule{\True}{\Diamond}[2]}
                [\pFmla{\False}{\formula{A}}{3},
                  just = {\TRule{\False}{\Diamond}[4]}, close]
            ]
          ]
        ]
      ]
    ]
  \end{oltableau}
\end{ex}

\end{document}